\scanpage{47}
\mspage{134}

comme dans le lemme :
\diagram{\begin{tikzcd}[ampersand replacement=\&,row sep=3pt,column sep=3pt]
X'\arrow[rr,"f"]\arrow[ddr,"p'"']\&\&X\arrow[ddl,"p"]\&\&\&\\
\&\&\&\overline{X'}\arrow[rr,"\bar f"]\arrow[ddr,"\overline{p'}"']\&\&\bar X\arrow[ddl,"\bar p"]\\
\&S\&\&\&\&\\
\&\&\&\&\bar S\arrow[ulll,"\unsure{\varphi|q}"']\&
\end{tikzcd}}.

Alors $\overline{p'}$ et $\bar p$ sont des Hot-fibrations, et
comme $\varphi$ est une Hot-fibration, le chgt. de
base par $\varphi$ transforme la Hot-équiv. $f$ en
une Hot-équiv. $\bar f$. Considérons alors
$\bar s\in\mathrm{Ob}\,\bar S$, et le diagramme
\diagram{\begin{tikzcd}[ampersand replacement=\&,row sep=small,column sep=small]
\overline{X'}\arrow[r,"\bar f"]\&\bar X\\
\overline{X'}_{\bar s}\arrow[u,hook,"i'"]\arrow[r,"\bar f_{\bar s}"]\&\bar X_{\bar s}\arrow[u,hook,"i"']
\end{tikzcd}}

Comme $\{\bar s\}\longrightarrow\bar S$ est une Hot-équivalence,
la flèche $i$ qui s'en déduit par le chgt.
de base $\bar p:\bar X\to\bar S$ l'est aussi puisque $\bar p$
est une Hot-fibration, et de $=$ pour
$i'$. Ainsi dans le carré \inserted{comm.} précédent, $\bar f$,
$\unsure{i|j}$ et $i'$ sont des Hot-équivalences, donc aussi
$\bar f_{\bar s}$. Mais \deleted{\unsure{comme}} si $s=\varphi(\bar s)$, $\bar f_{\bar s}$ s'identifie
à $\deleted{\unsure{\bar f_s}}f_s:X'_s\longrightarrow X_s$. D'ailleurs, il est
immédiat que \deleted{\unsure{si $S$}} $\varphi$, étant une Hot-fibration

% archive sheet number: 47
